% Generated draft: requires mathematical and source review.
\documentclass[11pt]{article}
\usepackage[margin=25mm]{geometry}
\usepackage{fontspec}
\setmainfont{Latin Modern Roman}
\usepackage{amsmath,amssymb,mathtools,unicode-math}
\setmathfont{Latin Modern Math}
\usepackage{enumitem,needspace,xurl,hyperref}
\setlength{\parindent}{0pt}
\setlength{\parskip}{5pt}
\setlength{\emergencystretch}{3em}
\newcommand{\N}{\mathbb{N}}
\newcommand{\Z}{\mathbb{Z}}
\newcommand{\Q}{\mathbb{Q}}
\newcommand{\R}{\mathbb{R}}
\newcommand{\C}{\mathbb{C}}
\newcommand{\F}{\mathbb{F}}
\newcommand{\A}{\mathbb{A}}
\newcommand{\PP}{\mathbb{P}}
\newcommand{\E}{\mathbb{E}}
\newcommand{\eps}{\varepsilon}
\newcommand{\dd}{\,\mathrm{d}}
\begin{document}
\section*{UnsolvedMath ID: 30004424}
\textbf{Source problem code:} catalogue entry
\subsection{Definitions and mathematical statement}
For each integer \(\ell\ge 2\), define the partitions
\[
\lambda_\ell=(3\ell-3,\underbrace{2,\ldots,2}_{\ell-1\text{ entries}}),\qquad
\mu_\ell=(\ell+1,\underbrace{1,\ldots,1}_{\ell-2\text{ entries}}).
\]
For \(\ell=2\), the second repeated list is empty, so \(\mu_2=(3)\). For a vector space \(E\) over a specified field, \(\operatorname{Sym}^{d}E\) denotes its \(d\)-th symmetric power. The notation \(\nabla_\lambda\) is retained from the source for its partition-indexed representation construction; its precise convention is not supplied.

The source asks for a geometric or invariant-theoretic interpretation, for every integer \(\ell\ge2\), of the displayed isomorphism
\[
\nabla_{\lambda_\ell}\bigl(\operatorname{Sym}^{\ell}E\bigr)
\cong
\nabla_{\mu_\ell}\bigl(\operatorname{Sym}^{3\ell-4}E\bigr).
\]
This is a request to explain the supplied isomorphism, rather than a separately specified theorem to prove. The record does not identify the ground field or its characteristic, the dimension of \(E\), the acting group, the category in which the isomorphism holds, or the convention for \(\nabla_\lambda\). It also gives no criterion for what constitutes the requested interpretation. Consequently, the display is preserved as a source assertion with unresolved hypotheses; it is not asserted here for arbitrary vector spaces or representation categories.
\subsection{Short English statement}
For every integer \(\ell\ge2\), seek a geometric or invariant-theoretic explanation of the stated isomorphism between two partition-indexed constructions on symmetric powers of \(E\). The setting of the representations and the required form of the explanation are unspecified and require clarification.
\subsection{Sources}
\begin{itemize}
\item Local UnsolvedMath registry, \texttt{data/unsolvedmath/problems.json}, permanent ID \(30004424\), source problem code \texttt{catalogue entry}, title: Geometry behind a Plethystic Representation Isomorphism. The supplied record belongs to the Oberwolfach Reports Open Problems collection. No individual proposer is identified; the supplied proposed year is 2020.
\item Supplied source citation: Mini-Workshop: Kronecker, Plethysm, and Sylow Branching Coefficients and their Applications to Complexity Theory (2021), Oberwolfach Reports, DOI 10.4171/owr/2020/9, pp.\ 517--568. Source URL: \url{https://doi.org/10.4171/owr/2020/9}. The background describes these reports as documenting mathematical workshops and their open-problem discussions.
\item Supplied literature source, labelled ``original OWR report'': \url{https://doi.org/10.4171/owr/2020/9}.
\item Literature assessment checked at \texttt{2026-08-22}: the supplied curation labels the problem open because the requested geometric or invariant-theoretic interpretation has no precise target theorem and remains unresolved or unclear. This is supplied dated curation, not independently verified literature assessment. The source explicitly requests independent verification of time-sensitive status and cited claims.
\item The registry labels the statement \texttt{corrected\_verified} and reports a manual semantic, typographic, and LaTeX review claiming that the clean statement is self-contained and preserves its target without unsupported additions. That supplied review label does not resolve the missing representation-theoretic hypotheses identified here.
\end{itemize}
\end{document}
